\documentclass[tikz,border=10pt]{standalone}
\usepackage{tikz}
\usepackage{amsmath}
\usepackage{amssymb}
\usepackage{ctex}
\usetikzlibrary{calc, positioning, arrows.meta, decorations.pathreplacing, patterns, shadings}

\begin{document}
\begin{tikzpicture}[>=Stealth, scale=1.1]

% ================================================================
% 列空间平面 C(A)（浅蓝色平行四边形表示子空间）
% ================================================================
\fill[blue!10, draw=blue!40, rounded corners=2pt]
  (-2.0, -1.0) -- (3.5, -0.5) -- (2.5, 1.5) -- (-1.0, 1.0) -- cycle;

\node[font=\small, blue!70!black] at (0.8, 1.2) {列空间 $C(A)$};

% ================================================================
% 原点 O（在列空间内）
% ================================================================
\coordinate (O) at (0, 0);
\fill (O) circle (1.5pt) node[below left, font=\scriptsize] {$O$};

% ================================================================
% 投影点 b_hat（在列空间内）
% ================================================================
\coordinate (BHAT) at (1.6, 0.4);
\fill[red!80!black] (BHAT) circle (2pt);
\node[font=\small, below right, red!80!black] at (BHAT) {$\hat{\mathbf{b}} = A\hat{\mathbf{x}}$};

% ================================================================
% 目标点 b（在列空间外的 R^m 中）
% ================================================================
\coordinate (B) at (1.6, 2.8);
\fill[purple!80!black] (B) circle (2pt);
\node[font=\small, right, purple!80!black] at (B) {$\mathbf{b}$};

% ================================================================
% 残差向量 r = b - b_hat（垂直于列空间）
% ================================================================
\draw[->, very thick, green!60!black] (BHAT) -- (B)
  node[midway, right, font=\small, green!60!black] {$\mathbf{r} = \mathbf{b} - \hat{\mathbf{b}}$};

% 直角符号（r ⊥ C(A)）
\draw[gray!60, thin]
  ($(BHAT)+(0.18, 0)$) -- ($(BHAT)+(0.18, 0.18)$) -- ($(BHAT)+(0, 0.18)$);

% ================================================================
% 垂直虚线（从 b 到 b_hat）
% ================================================================
\draw[densely dashed, gray!50] (B) -- (BHAT);

% ================================================================
% 列空间内的向量（两个基向量，表示 A 的列）
% ================================================================
\draw[->, thick, blue!60!black] (O) -- ($(O)+(2.0,-0.2)$)
  node[below right, font=\scriptsize] {$\mathbf{a}_1$};
\draw[->, thick, teal!60!black] (O) -- ($(O)+(0.8,0.9)$)
  node[above right, font=\scriptsize] {$\mathbf{a}_2$};

% ================================================================
% b_hat 从 O 出发的向量
% ================================================================
\draw[->, thick, red!70!black] (O) -- (BHAT);

% ================================================================
% "最短距离"标注
% ================================================================
\draw[<->, gray!60, thin]
  ($(B)+(-0.15, 0)$) -- ($(BHAT)+(-0.15, 0)$)
  node[midway, left, font=\scriptsize, gray!70] {最短距离 $\|\mathbf{r}\|$};

% ================================================================
% "正交条件"标注（A^T r = 0）
% ================================================================
\node[font=\scriptsize, fill=green!8, draw=green!50, rounded corners=2pt,
      inner sep=4pt, align=center]
  at (3.8, 2.5) {
    正交条件\\
    $A^T\mathbf{r} = \mathbf{0}$\\
    $\mathbf{r} \perp C(A)$
  };

% 指向直角的箭头
\draw[->, gray!50, thin] (3.5, 2.0) -- (1.9, 0.6);

% ================================================================
% 公式框（正规方程）— 放在图外下方
% ================================================================
\node[font=\small, fill=white, draw=gray!50, thin, rounded corners=3pt,
      inner sep=8pt, align=center]
  at (1.2, -2.0) {
    正规方程（Normal Equation）：
    $A^T A\,\hat{\mathbf{x}} = A^T\mathbf{b}$\\
    $\hat{\mathbf{x}} = (A^TA)^{-1}A^T\mathbf{b}$，\quad
    $\hat{\mathbf{b}} = A(A^TA)^{-1}A^T\mathbf{b}$
  };

% ================================================================
% 标题
% ================================================================
\node[font=\large\bfseries] at (1.2, 3.7) {最小二乘法的几何意义};

\end{tikzpicture}
\end{document}
